% !TeX root = ../../../main.tex

\section{Discussion}
\subsection{Interpreting Individual Figures}

Fig.~\ref{fig:boxplotRe} shows that VisR RE were significantly different for porcine kidneys with acute, chronic, and combined IRIs compared to kidneys without any injury. This distinction was observed in the outer and inner parts of the renal cortex, as well as for the outer and inner medulla, for all injury types except for the outer cortex of the combined injury cohort. VisR RE values of uninjured kidneys appeared to be higher in the inner cortex and outer medulla compared to the outer cortex and inner medulla. However, RE values generally appeared to drop for injured kidneys from all cohorts with respect to uninjured kidneys, suggesting a softening of shear elasticity along the direction of nephrons due to IRI. This drop in shear elasticity is particularly pronounced in the acute injury cohort in the medulla, as well as all cohorts in the inner cortex.

While this drop in RE initially appears to be counterintuitive (especially for the chronic injury cohort), Fig.~\ref{fig:boxplotRv} shows that RV also progressively increased from uninjured kidneys to those with acute, chronic, and combined injuries. Furthermore, differences in RV were most prominent in the outer and inner medulla, where values from the combined injury cohort were particularly large compared to the other cohorts. This is consistent with expectations that acute injuries involved interstitial edema and while chronic injuries involved fibrosis (and thereby changes in perfusion), both of which may manifest as increased viscosity in the medulla.

Similarly, Fig.~\ref{fig:boxplotDopio} demonstrates that DoPIo modulus estimates were distributed in a significantly different manner for porcine kidneys with acute, chronic, and combined IRIs compared to kidneys without any injury. This distinction was observed in the outer and inner parts of the renal cortex, as well as for the outer and inner medulla, for all injury types except for the outer cortex in chronic IRI. In particular, DoPIo-based stiffness estimates were also generally softer in injured kidneys versus controls, and the difference in the two distributions became more pronounced in structures closer to the pelvis as well as in the acute and combined injury cohorts.

Recall that DoPIo modulus values are presumably quantitative as they are calculated using a linear model based on times-intersect empirically observed in finite element model simulations. This stands in contrast to VisR RE values since they are qualitative biomarkers; RE correlates with both the elasticity of the interrogated tissue as well as the applied ARF force amplitude of unknown magnitude. Still, significant differences were found between both VisR RE and DoPIo modulus estimates in the uninjured renal parenchyma versus those from injured kidneys. This suggests that DoPIo and VisR are similarly capable of parameterizing the elasticity of anisotropic tissues \emph{in vivo} in a controlled imaging situation. Furthermore, both techniques can be used to distinguish renal parenchyma with acute, chronic, and combined injury cohorts from uninjured parenchyma in various structural regions.

However, there were two exceptions to how DoPIo and VisR agreed in which anatomical regions indicated significantly different mechanical properties in ARFI-displaced tissue due to pathology. First, in Fig.~\ref{fig:boxplotDopio}, DoPIo modulus estimates were not found to be significantly different between the outer cortices of uninjured kidneys versus those from the chronic injury cohort. This stands in contrast to Fig.~\ref{fig:boxplotRe}, where RE values at the same anatomical feature were found to be significantly different between uninjured and all three injured cohorts except for the combined injury cohort.

The ability to distinguish between outer versus inner regions of cortices and medullae using DoPIo and VisR is significant in the context of renal pathologies. The outer cortex is predominantly populated by glomeruli and interstitia while the inner cortex contains more tubules, the structures where increased oxygen consumption occurs. As a result, certain insults such as oxygen deprivation are more impactful in the inner cortex, while others such as glomerulosclerosis may be present in higher densities in the outer cortex due to the lack of tubules. This difference is more extreme in the medulla where glomeruli are not present at all and oxygen demand is even higher. This expectation of higher ischemic injury in the cortex is consistent with Fig.~\ref{fig:boxplotDopio} since elasticity estimates in deeper kidney regions, such as those in the inner medulla, are generally lower than in more superficial regions like the outer cortex in injured cohorts.

Unlike Figures \ref{fig:boxplotRe} and \ref{fig:boxplotDopio}, VisR and DoPIo DoA measurements in Figures \ref{fig:boxplotReDoa}, \ref{fig:boxplotRvDoa}, and \ref{fig:boxplotDopioDoa} respectively did not show as many significant differences. Fig.~\ref{fig:boxplotReDoa} only indicated significantly different DoA values for the outer cortex in the acute injury cohorts versus controls, as well as for all three injury cohorts versus controls in the inner cortex. However, the differences in RE-based ln(DoA) in the inner cortex were due to the range of values rather than due to any median shifts. However, RV-based ln(DoA) were only significantly increased in the combined injury cohort's outer cortex, but their distributions were also much wider in the outer and inner medulla. This suggests that the combined presence of acute and chronic IRI injuries greatly varied medullary viscous anisotropy, though more samples are required to make any inferences about their directionality. On the other hand, DoPIo indicated significantly different DoA values in kidneys with chronic and combined injuries versus controls in the outer medulla - and this difference involved a modest decrease in median DoA values for injured kidneys. Since the medulla is more poorly perfused than the cortex and IRI is expected to lead to collagen deposition in the interstitium between nephrons, these changes in the outer medulla are consistent with expectations. However, the veracity of shear modulus estimates in this context, as well as explanations for why VisR RE and DoPIo disagree in which anatomical regions show significant differences as well as why the DoPIo results do not extend to the inner medulla, require further study.

RR measurements for VisR and DoPIo were significantly different versus uninjured kidneys, indicating differences in elasticity between structures in the renal parenchyma due to pathology. In Fig.~\ref{fig:boxplotReRr} and \ref{fig:boxplotRvRr}, VisR RE- and RV-derived RR were significantly different between uninjured kidneys and the acutely injured kidneys for all six possible combinations of parenchymal regions except between outer and inner medulla for RV. Likewise, both RE- and RV-based RR for injured kidneys in the chronically injured cohort were found to be significantly different in the ratio of inner cortex to outer medulla. For kidneys that underwent combined injuries, the ratio of RE in the outer cortex versus inner cortex as well as outer cortex to inner medulla were significantly different versus those from control kidneys. This trend is consistent with the expectation that the structure of renal parenchyma is altered by ischemic injury, and that changes led to the cortex appearing softer than the medulla. Interestingly, RV-based RR for the combined injury cohort were significantly different from controls in ratios involving the inner cortex and outer medulla, a different pattern from those observed by RE. 

DoPIo ln(RR) measurements in Fig.~\ref{fig:boxplotDopioRr} were only significantly different between uninjured kidneys and the acutely injured kidneys for the comparison of the outer cortex to outer medulla. Likewise, chronically injured kidneys were found to be significantly different from uninjured kidneys for RR between outer and inner cortices. However, injured kidneys in the combined injury cohort were found to be significantly different from uninjured kidneys for RR between the outer cortex and outer medulla, inner cortex~/ outer medulla, and outer~/ inner cortices. This trend suggests that, though to a lesser degree than VisR, DoPIo is also capable of detecting relative differences in elasticity between regions of the renal parenchyma in the presence of pathology. The two techniques mainly agreed in cases involving the outer cortex.

\subsection{Detecting Renal Pathology with VisR and DoPIo}

Fig.~\ref{fig:auc} shows the performances of binary classifiers developed using VisR RE and DoPIo modulus measurements. Starting with the top left bars in Fig.~\ref{fig:auc}(a) to compare classifiers using VisR RE values alone, RV values alone, and those using DoPIo modulus estimates alone, AUC distributions were significantly different between RE and DoPIo. This suggests that DoPIo modulus estimates were better able to detect kidneys with any pathology than VisR RE values. However, RV, which fell in between the two for median and MAD AUC, were not significantly different from either. Proceeding down the remaining groups of bars, for a given classification problem, classifiers that used multiple parameters (e.g. models using raw values and RR) did not substantially outperform those that solely relied on raw values. This suggests that a single VisR or DoPIo acquisition of a single part of the renal parenchyma may be sufficient for detecting the presence of renal pathologies.

Continuing down the rest of the "Any Injuries" subfigure: classifiers that solely considered RR inputs, as well as those that involved RE/modulus inputs in combination with RR and/or DoA, observed minimal differences in AUC distributions when DoPIo modulus estimates were inputted rather than VisR RE or RV. One exception of this is for RR where RV outperformed against RE, though neither significantly outperformed DoPIo. This suggests that RV offers marginal improvements in identifying kidneys with any injury using RR, but DoPIo is superior over VisR in detecting renal parenchymal injuries using single measurements. The median and MAD of AUC values did not noticeably change when RR and/or DoA inputs were added to RE/modulus inputs for both VisR- and DoPIo-based models, though DoPIo outperformed both RE and RV when raw parameters, RR, and DoA were all allowed as model inputs. However, the resulting AUC for DoPIo was nearly identical to median classifiers that solely used raw values. This suggests that, while RR and DoA inputs may provide additional information in detecting renal pathologies, they do not add significant diagnostic value when combined with RE or modulus information.

Looking at the remaining subfigures of Fig~\ref{fig:auc}, DoPIo modulus-based models also outperformed VisR RE-based models for distinguishing kidneys with combined acute-chronic injuries from all others. While differences in AUC distributions were significant using RE/modulus inputs alone, the MAD of AUC slightly increased when classifiers also allowed RR or both RR and DoA to be included in inputs. This suggests that, as with the task of detecting any injury, DoPIo modulus estimates, alone, offer improved differentiation of kidneys with simultaneous acute and chronic injuries versus RE. RR and DoA, as well as RV, offer additional information, but do not add significant diagnostic value when combined with modulus information.

However, this comparison changes when considering the tasks of distinguishing kidneys with acute injuries or chronic injuries from all others. Numerically, AUC values for classifiers of acutely-injured kidneys were often higher for VisR-based metrics than for DoPio. However, many of these differences were not statistically significant. The only exception is for classifiers solely using DoA where, uniquely for this classification task, RV-based detectors significantly outperformed the RE-based ones. While this suggests that RE-based DoA was a uniquely poor standalone biomarker for detecting acute injuries, the addition of DoA to RE/RV or RR neither helped nor hindered the performance of classifiers. This may be due to how, in Fig.~\ref{fig:boxplotRe}-\ref{fig:boxplotDopioRr}, there were few comparisons where acutely injured kidneys significantly "stood out" both from controls and other injured kidneys in a unique way.

For the task of detecting kidneys with only chronic injuries, classifiers that used VisR RE- and RV values often showed higher AUC compared to DoPIo modulus-based ones. However, the only case where the AUC distributions were significantly different was for RR between RV and DoPIo. Still, this potentially suggests that models that considered RE/RV inputs (on its own, as well as when combined with either RR or DoA inputs), may be more performant than those based on DoPIo. This would suggest, then, that VisR RE and/or RV values may be more sensitive to the presence of chronic injuries in the renal parenchyma than DoPIo modulus estimates. However, this observation is tentative since most differences in AUC distributions were not statistically significant.

\subsection{Histological Validation}

Differences in VisR and DoPIo elasticity metrics also correlated with histological differences in Fig.~\ref{fig:histo} and Table~\ref{tbl:histoQuant}. Histological changes were consistent with IRI as expected in the three injured cohorts; the acutely injured kidneys showed immune responses to the induced injury, the chronic injury cohort showed cortical fibrosis, and the combined injury cohort showed aspects of both acute and chronic injuries.

In the acute IRI group, kidneys had marked widespread renal cortical tubular epithelial degeneration, regeneration, and necrosis. In the renal cortex, some tubules had epithelial loss and attenuation that suggests necrosis. Tubules also contained intratubular proteinaceous material, desquamated tubular epithelial cells, and macrophages. Regions of cortex also had robust tubular epithelial regenerative changes. There was often minimal or mild interstitial edema in the cortex and medulla. By IHC, there were notable differences in Iba1-positive inflammatory cell infiltrates in the cortex of IRI kidneys (mostly of the "marked" score) compared to uninjured kidneys ("minimal" and "mild"). Neutrophils were occasionally evident within the cortical and medullary tubules of IRI kidneys. Multifocal tubular mineralization, a common sign of cell death, was more prominent in the cortex than medulla. Fibrosis related to the IRI procedure was not apparent in any sections by H\&E or Sirius Red histochemical staining.

In the chronic IRI group, kidney injuries were associated with mild to moderate segmental cortical interstitial fibrosis, mostly along nephron segments. The injuries were evident by picro-Sirius red histochemical staining in the chronically injured kidneys, but they were minimal in the control kidneys as expected. Within the fibrotic regions were mostly mild infiltrates of Iba1-positive inflammatory cells (by IHC) compared to minimal infiltrates in control kidneys. Scattered glomeruli were sclerotic (fibrosed/nonfunctional) and there was some tubular dropout (tubules injured and lost). There is no notable difference between control and IRI kidneys in the medulla.

In the combined cohort, acute findings were identical to those in the acute IRI group, and chronic findings were similar to the chronic IRI group. Interestingly, the chronic changes (i.e. fibrosis, tubular loss/dropout) were slightly more extensive in this group than in the chronic IRI only group. In this group, the severity of the interstitial fibrosis was mostly moderate to marked (rather than mild to moderate as in the chronic only IRI group). Iba1-positive inflammatory cell infiltrates (by IHC) were at a similar incidence and severity compared to the acute IRI kidneys.

The agreement of histological changes to pig cohort assignments in Table~\ref{tbl:histoQuant} show that elastic changes observed using VisR and DoPIo are consistent with the presence of renal pathologies. The lesions found in the injured kidneys of the acute injury cohort include inflammatory cell infiltrates (i.e. inflammation) in the cortex and medulla, whereas those in the chronic injury cohort presented with a significantly different amount of cortical fibrosis compared to the uninjured controls. Both types of lesions were present to a statistically significant degree in the combined injury cohort. Considering that the presence of both inflammation and fibrosis is a hallmark of AMR while many related pathologies only have one or the other~\cite{loupy_Banff_2020}, these results suggest that both VisR and DoPIo noninvasively detect histological changes that are consistent with AMR in transplanted kidneys.

It should be noted, however, that some pathological findings were also present in contralateral kidneys that did not undergo IRI. Mild interstitial edema detected in the uninjured kidneys of the combined injury cohort, as well as macrophages and fibrosis in all three cohorts' contralateral kidneys, were present unlike in the uninjured controls. This suggests that systemic responses to IRI took place that affect both kidneys.

\subsection{On Limitations of This Study}

As evident in Figures \ref{fig:boxplotRe}, \ref{fig:boxplotRv}, and \ref{fig:boxplotDopio}, this study treated VisR RE, RV, and DoPIo modulus estimates as directly comparable metrics. This interpretation was possible since the kidney was exteriorized during imaging, allowing focusing of ARF excitation and tracking beams to be highly standardized. In clinically relevant \emph{in vivo} imaging settings, kidneys would be imaged with abdominal tissue in between the transducer and the kidney. The thickness, composition, and geometry of this tissue would vary by patient, as well as within the same patient over time, causing the acoustic energy distribution of ARF excitations to vary across patients and over time. Therefore, the direct comparison of VisR RE and DoPIo modulus estimates is not likely to be relevant in more realistic imaging scenarios. In such cases, only DoPIo modulus estimates from clinical acquisitions may be compared to this study; VisR RE, as a qualitative parameter, would not be directly comparable. Still, since the feasibility of on-axis ARF-based imaging has been demonstrated previously~\cite{hossain_Evaluating_2018} following an earlier version of the animal model used in this study~\cite{hossain_Mechanical_2019}, the results of this study nevertheless provide a useful benchmark for future \emph{in vivo} imaging studies in clinically relevant environments.

Individual imaging situations also presented challenges in this study. Since the pigs imaged in this study were adolescent, some kidneys were larger than others. This led to discrepancies in which kidneys could be directly imaged by the transducer, and which kidneys required the use of a standoff pad given the fixed focal depth of 25 mm. This specific focal depth was required since empirical models for converting times-intersect to modulus estimates for DoPIo were developed specifically using a 25 mm focal depth for push and track beams. This approach allowed consistency to be maintained between elastic modulus estimates in pig kidneys and training data used to create the model. While this approach allowed all VisR and DoPIo images from all cohorts to be compared in a standardized way, this inconsistent use of standoff pads may have introduced variability in both DoPIo modulus estimates and VisR RE values.

While this study was designed to assess whether VisR and DoPIo can differentiate kidneys with inflammation and/or fibrosis from each other and from uninjured kidneys, the resulting RE, RV, and estimated shear modulus values for the parenchyma were not directly measured using a validation standard. This was because the kidneys were imaged \emph{in vivo} and were not excised until after imaging was complete, and shear wave imaging was not performed alongside VisR and DoPIo. Recent review articles have shown a large degree of variability in shear wave speeds in renal parenchyma - including discrepancies in expected shear modulus values for the medulla versus cortex - due to differences in imaging plane orientation, renal perfusion and blood flow, glomerular filtration rate (i.e. renal function), whether the kidney is the dominant one in the body, and whether multifocal lesions were sufficiently included in the imaging plane~\cite{urban_Novel_2021,lim_Shear_2021,ce_Ultrasound_2023,leong_Stiffness_2020}. Future clinical work should control those factors such that a more appropriate baseline can be estabished for elastic parameters of renal parenchyma. This would also allow for better validations of VisR and DoPIo's performance against off-axis elastography techniques.

The DoPIo empirical models used in this study were developed using elastically isotropic materials~\cite{yokoyama_DoubleProfile_2025}. However, the renal parenchyma is more appropriately considered an anisotropic medium. Muhtadi \emph{et al.}~\cite{muhtadi_Double_2025} have shown that, when DoPIo is applied to \emph{in silico} TI materials such that an ARF push is applied perpendicular to the material's axis of symmetry (AoS) and the pair of tracking beams have different widths in the direction of the AoS, the longitudinal shear modulus is primarily estimated. However, this modulus estimation can be confounded by two factors: (1) if the ARF push is not perfectly perpendicular to the AoS, and (2) if the longitudinal Young's modulus is increased. These confounding factors were also reported for VisR~\cite{muhtadi_VisR_2024} as well as in shear wave imaging~\cite{andrade_Simultaneous_2025,ngo_plane_2024}. While we standardized ROI positions based on the presumed orientation of the parenchyma's axis of symmetry to control for this variability, discrepancies in nephron orientation may cause the first confounding factor to manifest. Furthermore, IRI could cause changes in the transverse shear elastic moduli and longitudinal Young's moduli of renal parenchyma, as well, manifesting in the second confounding factor. The specific quantitative impacts of changes in all four elastic moduli of TI materials on DoPIo's performance remain an open question to be tackled in a similar manner as done for VisR~\cite{muhtadi_VisR_2024}.

Similarly, while the renal parenchyma is viscoelastic, DoPIo is currently developed with the assumption of purely elastic media. Since viscous materials can delay displacement recovery behavior, DoPIo times-intersect could be delayed in viscoelastic materials. This may cause DoPIo's empirically derived model, which was only trained on simulated media that were elastic and isotropic, to underestimate the shear elasticity of more viscous materials.  These combined issues of viscosity and anisotropy, as well as differences in how those effects manifest between DoPIo and VisR, could explain their disagreements in hypothesis test results. However, the development of a framework to systematically characterize the impacts of viscosity and anisotropy to DoPIo are beyond the scope of this study and topics of ongoing work.

Finally, the animal model used in this study reflects structural changes in the renal parenchyma seen in AMR, but it is not a direct model of AMR itself. Inflammation and/or fibrosis was induced through vascular occlusion and reperfusion in a native kidney, rather than through immune responses to transplanted organs. The IRI protocol also applied injuries throughout the entire kidney, whereas some early-stage lesions associated with transplant rejections may be more focal in nature~\cite{saritas_Kidney_2021,loupy_Banff_2020}. Nevertheless, the presence of both inflammation and fibrosis in the injured cohorts is consistent with histological findings in AMR~\cite{loupy_Banff_2020,venner_Relationships_2016}, suggesting that VisR and DoPIo may be useful for detecting renal transplant pathologies. Future work will involve imaging transplanted kidneys in animal models that more closely resemble human renal transplant rejection.

